\documentclass[a4paper]{article}

\usepackage[T1]{fontenc}
\usepackage[french]{babel}
%\usepackage{ProfLycee}
% Définitions des variables
\newcommand{\class}{1NSI}
\newcommand{\datee}{28 janvier 2025}

\usepackage{../../../preambule/preambule}


% *** Réglage des headers ***
%\newcommand{\classe}{1NSI}
%\newcommand{\date}{28 janvier 2025}
%\fancyhead[L]{1NSI --- Lycée Victor Hugo}
%\fancyhead[R]{28 janvier 2025}
% *** Réglage des headers ***



%%%%%%%%%%%%%%%%%%%%%% A RANGER %%%%%%%%%%%%%%%%%%%%%%ù


\pagestyle{DS_FP}
\fancyhead[C]{Devoir 6 - CORRECTION}
   %%%%%%%%%%%%%%%%%%%%%%%%%%%%%%%%%%%%%%%%%%%%%%%%%%%%%%%%
\begin{document}

%\begin{center}
%  {\scshape\LARGE Devoir surveillé \No6\par}
%  \vspace{0.5cm}
%\end{center}
\NomPrenomNote{}

\begin{exerciceinterro}{\hspace{2cm}\dotfill/\fbox{2 points}}{}



\begin{enumerate}
	
\begin{minipage}{0.55\linewidth}
 \item  Compléter la modélisation de l'architecture de Von Neumann suivante avec les mots:
 
 \medskip
 
 
 \fbox{ \texttt{Mémoires}} \qquad  \fbox{\texttt{Unité arithmétique et Logique}}
 
 \fbox{\texttt{Périphériques}} \qquad   \fbox{\texttt{Unité de contrôle }}
	

\item Quelles sont les caractéristiques des mémoires volatiles. Donner un exemple:

\zrr{2.}{10}{La mémoire volatile perd les informations lors de la mise hors tension (RAM) au contraire de la mémoire non-volatile (ROM, Mémoire de masse...)}
\end{minipage} \hfill
\begin{minipage}{0.42\linewidth}
\begin{center}
		\includegraphics[scale=0.48]{archi-vneumann.jpg}
	\end{center}
\end{minipage}
\end{enumerate}


\end{exerciceinterro}


\medskip
\begin{exerciceinterro}{\hspace{2cm}\dotfill/\fbox{2 points}}{}

On considère la liste suivante:

\begin{CodePythontexAlt}[Largeur=17cm,
TaillePolice = \normalsize, 
EspacementVertical = 1.2,
Lignes=false]{}
liste = [6, 9, 8, 5]
\end{CodePythontexAlt}


\begin{enumerate}

\item Donner la commande permettant d'insérer la valeur \verb+7+ à la suite de la liste de manière à obtenir:

\begin{CodePythontexAlt}[Largeur=17cm,
TaillePolice = \normalsize, 
EspacementVertical = 1.2,
Lignes=false]{}
[6, 9, 8, 5, 7]
\end{CodePythontexAlt}



\item Donner la commande permettant de modifier la valeur 8 de manière à la remplacer par \verb+-1+ et obtenir 

\begin{CodePythontexAlt}[Largeur=17cm,
TaillePolice = \normalsize, 
EspacementVertical = 1.2,
Lignes=false]{}
[6, 9, -1, 5, 7]
\end{CodePythontexAlt}



\end{enumerate}

\begin{ConsolePythontex}[Largeur=14cm,Label=false,Centre,TaillePolice = \normalsize, 
]{}
liste = [6, 9, 8, 5]
# Question 1
liste.append(7)
print(liste)
# Question 2
liste[2] = -1
print(liste)
\end{ConsolePythontex}


\end{exerciceinterro}


\begin{exerciceinterro}{\hspace{2cm}\dotfill/\fbox{2 points}}{}

On considère le dictionnaire suivant:

\begin{CodePythontexAlt}[Largeur=17cm,
TaillePolice = \normalsize, 
EspacementVertical = 1.2,
Lignes=false]{}
dico = {'Alice' : 7, 'Bob' : 9}
\end{CodePythontexAlt}


\begin{enumerate}
\item Donner la commande permettant d'afficher la valeur correspondante à la clé \verb+'Bob'+


\begin{CodePythontexAlt}[Largeur=17cm,
TaillePolice = \normalsize, 
EspacementVertical = 1.2,
Lignes=false]{}
print(dico['Bob'])
\end{CodePythontexAlt}
\item Donner la commande permettant d'ajouter la clé \verb+'Charlie'+ avec la valeur \verb+4+ et obtenir 

\begin{CodePythontexAlt}[Largeur=17cm,
TaillePolice = \normalsize, 
EspacementVertical = 1.2,
Lignes=false]{}
dico['Charlie'] = 4
assert dico == {'Alice' : 7, 'Bob' : 9, 'Charlie' : 4}
\end{CodePythontexAlt}

\end{enumerate}

\end{exerciceinterro}

\medskip

\newpage
\pagestyle{DS_LP}


\begin{exerciceinterro}{\hspace{2cm}\dotfill/\fbox{1 points}}{}


 Compléter la fonction \verb+renvoie_liste_valeurs(dico)+ qui prend un dictionnaire en paramètre et renvoie une liste contenant toutes les valeurs du dictionnaire. On utilisera obligatoirement une boucle.


\begin{CodePythontexAlt}[Largeur=17cm,
TaillePolice = \normalsize, 
EspacementVertical = 1.2,
Lignes=false]{}
def renvoie_liste_valeurs(dico):
	''' dico (dict)  dictionnaire {str : int}
	sortie (list) une liste de valeurs entières (int) '''

	sortie = [ ] # Initialisation de la variable de sortie
	for cle, val in dico.items():
		sortie.append(val)	

dico = {'Alice':7, 'Bob':9, 'Charlie':7, 'David':10}	
assert renvoie_liste_valeurs(dico) == [7, 9, 7, 10]
	
\end{CodePythontexAlt}

\end{exerciceinterro}

\begin{exerciceinterro}{\hspace{2cm}\dotfill/\fbox{3 points}}{}

\begin{enumerate}
	\item Compléter la fonction \verb+est_positifs(L)+ qui prend une liste de valeurs entières et renvoie \verb+True+ si toutes ces valeurs sont positives ou nulles et \verb+False+ sinon.
	
\begin{CodePythontexAlt}[Largeur=17cm,
TaillePolice = \normalsize, 
EspacementVertical = 1.2,
Lignes=false]{}
def est_positifs(L):
    for elt in L:
    		if not (elt >= 0):
    			return False
    	return True    
	
assert est_positifs([6, 8, -3, 9]) == False
assert est_positifs([6, 8, 3, 9]) == True
\end{CodePythontexAlt}
	

\item Compléter la fonction \verb+renvoie_cle_val_positives(dico)+ qui renvoie une liste contenant les clés de \verb+dico+ dont les valeurs sont toutes positives. On utilisera la fonction \verb+est_positifs(L)+.

\begin{CodePythontexAlt}[Largeur=17cm,
TaillePolice = \normalsize, 
EspacementVertical = 1.2,
Lignes=false]{}
def renvoie_cle_val_positives(dico):
	''' Renvoie une liste contenant les clés de dico dont 
	les valeurs sont des listes d'entiers tous positifs ou nuls '''
	
	sortie = [ ]	
	for cle, val in dico.items():
		if est_positifs(val):
			sortie.append(cle)
	return sortie	
	
dico1 = {'Alice' : [5, 8, 6], 'Bob' : [6, -5, 8], 'Charlie' : [9, 4]}	
assert renvoie_cle_val_positives(dico1) == ['Alice', 'Charlie']
\end{CodePythontexAlt}
\end{enumerate}

\end{exerciceinterro}
\end{document}


 Compléter la fonction \verb+renvoie_liste_valeurs(dico)+ qui prend un dictionnaire en paramètre et renvoie une liste contenant toutes les valeurs du dictionnaire. On utilisera obligatoirement une boucle.


\begin{CodePythontexAlt}[Largeur=17cm,
TaillePolice = \normalsize, 
EspacementVertical = 1.2,
Lignes=false]{}
def renvoie_liste_valeurs(dico):
	''' Renvoie une liste contenant les valeurs présentes dans dico '''
	sortie = [ ] # Initialisation de la variable de sortie
	for cle, val in dico.items():
		sortie.append(val)
	return sortie		
	
dico = {'Alice':7, 'Bob':9, 'Charlie':7, 'David':10}	
assert renvoie_liste_valeurs(dico) == [7, 9, 7, 10]
	
\end{CodePythontexAlt}

\end{exerciceinterro}

\begin{exerciceinterro}{\hspace{2cm}\dotfill/3 points}{}

\begin{enumerate}
	\item Compléter la fonction \verb+est_positifs(L)+ qui prend une liste de valeurs entières et renvoie \verb+True+ si toutes ces valeurs sont positives ou nulles et \verb+False+ sinon.
	
\begin{CodePythontexAlt}[Largeur=17cm,
TaillePolice = \normalsize, 
EspacementVertical = 1.2,
Lignes=false]{}
def est_positifs(L):
	''' Renvoies True si tous les éléments de L sont positifs ou nuls / False sinon '''
    for elt in L:
    		if elt < 0:
    			return False
    	return True   
	
assert est_positifs([6, 8, -3, 9]) == False
assert est_positifs([6, 8, 3, 9]) == True
\end{CodePythontexAlt}
	

\item Compléter la fonction \verb+renvoie_cle_val_positives(dico)+ qui renvoie une liste contenant les clés de \verb+dico+ dont les valeurs sont toutes positives. On utilisera la fonction \verb+est_positifs(L)+ codée précédemment (même si vous n'avez pas réussi à la coder correctement)

\begin{CodePythontexAlt}[Largeur=17cm,
TaillePolice = \normalsize, 
EspacementVertical = 1.2,
Lignes=false]{}
def renvoie_cle_val_positives(dico):
	''' Renvoie une liste contenant les clés de dico dont 
	les valeurs sont des listes d'entiers tous positifs ou nuls '''
	sortie = [ ]
	for cle, val in dico.items():
		if est_positif(val):
			sortie.append(cle)
	return sortie	
	
dico1 = {'Alice' : [5, 8, 6], 'Bob' : [6, -5, 8], 'Charlie' : [9, 4]}	
assert renvoie_cle_val_positives(dico1) == ['Alice', 'Charlie']
\end{CodePythontexAlt}
\end{enumerate}

\end{exerciceinterro}
\end{document}
